% GENERATED FILE. Edit canonical lessons, then run npm run generate:books.
% canonical-source-hash: fnv1a64:501c96367bd5dd41
% canonical-lessons: AR-C215-hawiyya, AR-C215-tashira, AR-C215-amtia, AR-C215-shahin, AR-C215-taksi

\chapter{Identity card, Visa, Luggage, Charger, Taxi}
\label{ch:ar-identity-card-visa-luggage-charger-taxi}
\hlchaptermodality{215}

\begin{chapteropening}
\textbf{By the end of this chapter:} \emph{I can say an identity card, a visa, luggage, a charger and a taxi.}

Complete the last lesson of chapter 215: I can say an identity card, a visa, luggage, a charger and a taxi.
\end{chapteropening}

\section[\texorpdfstring{\ar{هوية} (hawiyya)}{هوية (hawiyya)}]{\ar{هوية} (hawiyya) — an identity card}

\label{lesson:AR-C215-hawiyya}



\begin{quote}
Before the new one: say the Arabic for cotton, then the Arabic for wool.
\end{quote}

\subsection*{You'll want to know: \ar{هوية}}
\textbf{\ar{هوية}} — \emph{hawiyya} — \textquotedblleft{}an identity card\textquotedblright{}.

\begin{cousinweb}
\textbf{\ar{هوية}} (\emph{hawiyya}) is built from the pronoun \emph{huwa}, \textquotedblleft{}he\textquotedblright{}, with the abstract ending \emph{-iyya}: \textquotedblleft{}he-ness\textquotedblright{}, usually explained as a coinage of the medieval philosophers for identity.
\end{cousinweb}

\begin{grammarlens}[title={What you've built}]
One more word for this chapter.
\end{grammarlens}

\subsection*{Your turn}
\begin{itemize}
\raggedright
  \item \emph{Say it:} \emph{hawiyya}
  \item \emph{Say it:} \emph{hawiyya}, once more
  \item \emph{Say it:} say \emph{hawiyya}
  \item \emph{Recall:} say the Arabic for cotton, then the Arabic for wool, then say \emph{hawiyya} again
\end{itemize}

\begin{tcolorbox}[breakable,colback=teal!4,colframe=teal!35!black,title={Before you move on}]
What does \ar{هوية} mean? (\textquotedblleft{}An identity card\textquotedblright{}.) Say it once more.
\end{tcolorbox}
\section[\texorpdfstring{\ar{تأشيرة} (taʾshīra)}{تأشيرة (taʾshīra)}]{\ar{تأشيرة} (taʾshīra) — a visa}

\label{lesson:AR-C215-tashira}



\begin{quote}
Before the new one: say the Arabic for wool, then the Arabic for an identity card.
\end{quote}

\subsection*{You'll want to know: \ar{تأشيرة}}
\textbf{\ar{تأشيرة}} — \emph{taʾshīra} — \textquotedblleft{}a visa\textquotedblright{}.

\begin{cousinweb}
\textbf{\ar{تأشيرة}} (\emph{taʾshīra}) is built on the native Arabic root \emph{ʾ-sh-r}, \textquotedblleft{}to saw, to notch\textquotedblright{} (Lane). \emph{Taʾshīra} comes from the modern Form II verb \emph{ashshara}, \textquotedblleft{}to mark (a document)\textquotedblright{}, a sense not in Lane: a visa is a mark in a passport.
\end{cousinweb}

\begin{grammarlens}[title={What you've built}]
One more word for this chapter.
\end{grammarlens}

\subsection*{Your turn}
\begin{itemize}
\raggedright
  \item \emph{Say it:} \emph{taʾshīra}
  \item \emph{Say it:} \emph{taʾshīra}, once more
  \item \emph{Say it:} say \emph{taʾshīra}
  \item \emph{Recall:} say the Arabic for wool, then the Arabic for an identity card, then say \emph{taʾshīra} again
\end{itemize}

\begin{tcolorbox}[breakable,colback=teal!4,colframe=teal!35!black,title={Before you move on}]
What does \ar{تأشيرة} mean? (\textquotedblleft{}A visa\textquotedblright{}.) Say it once more.
\end{tcolorbox}
\section[\texorpdfstring{\ar{أمتعة} (amtiʿa)}{أمتعة (amtiʿa)}]{\ar{أمتعة} (amtiʿa) — luggage}

\label{lesson:AR-C215-amtia}



\begin{quote}
Before the new one: say the Arabic for an identity card, then the Arabic for a visa.
\end{quote}

\subsection*{You'll want to know: \ar{أمتعة}}
\textbf{\ar{أمتعة}} — \emph{amtiʿa} — \textquotedblleft{}luggage\textquotedblright{}.

\begin{cousinweb}
\textbf{\ar{أمتعة}} (\emph{amtiʿa}) is built on the native Arabic root \emph{m-t-ʿ}, \textquotedblleft{}to enjoy; goods\textquotedblright{} (Lane). \emph{Amtiʿa} is the plural of \emph{matāʿ}, \textquotedblleft{}goods, belongings\textquotedblright{}.
\end{cousinweb}

\begin{grammarlens}[title={What you've built}]
One more word for this chapter.
\end{grammarlens}

\subsection*{Your turn}
\begin{itemize}
\raggedright
  \item \emph{Say it:} \emph{amtiʿa}
  \item \emph{Say it:} \emph{amtiʿa}, once more
  \item \emph{Say it:} say \emph{amtiʿa}
  \item \emph{Recall:} say the Arabic for an identity card, then the Arabic for a visa, then say \emph{amtiʿa} again
\end{itemize}

\begin{tcolorbox}[breakable,colback=teal!4,colframe=teal!35!black,title={Before you move on}]
What does \ar{أمتعة} mean? (\textquotedblleft{}Luggage\textquotedblright{}.) Say it once more.
\end{tcolorbox}
\section[\texorpdfstring{\ar{شاحن} (shāḥin)}{شاحن (shāḥin)}]{\ar{شاحن} (shāḥin) — a charger}

\label{lesson:AR-C215-shahin}



\begin{quote}
Before the new one: say the Arabic for a visa, then the Arabic for luggage.
\end{quote}

\subsection*{You'll want to know: \ar{شاحن}}
\textbf{\ar{شاحن}} — \emph{shāḥin} — \textquotedblleft{}a charger\textquotedblright{}.

\begin{cousinweb}
\textbf{\ar{شاحن}} (\emph{shāḥin}) is built on the native Arabic root \emph{sh-ḥ-n}, \textquotedblleft{}to load, to fill\textquotedblright{} (Lane). \emph{Shāḥin} is the active participle, \textquotedblleft{}one that loads\textquotedblright{}; the charger is a modern sense.
\end{cousinweb}

\begin{grammarlens}[title={What you've built}]
One more word for this chapter.
\end{grammarlens}

\subsection*{Your turn}
\begin{itemize}
\raggedright
  \item \emph{Say it:} \emph{shāḥin}
  \item \emph{Say it:} \emph{shāḥin}, once more
  \item \emph{Say it:} say \emph{shāḥin}
  \item \emph{Recall:} say the Arabic for a visa, then the Arabic for luggage, then say \emph{shāḥin} again
\end{itemize}

\begin{tcolorbox}[breakable,colback=teal!4,colframe=teal!35!black,title={Before you move on}]
What does \ar{شاحن} mean? (\textquotedblleft{}A charger\textquotedblright{}.) Say it once more.
\end{tcolorbox}
\section[\texorpdfstring{\ar{تاكسي} (tāksī)}{تاكسي (tāksī)}]{\ar{تاكسي} (tāksī) — a taxi}

\label{lesson:AR-C215-taksi}



\begin{quote}
Before the new one: say the Arabic for luggage, then the Arabic for a charger.
\end{quote}

\subsection*{You'll want to know: \ar{تاكسي}}
\textbf{\ar{تاكسي}} — \emph{tāksī} — \textquotedblleft{}a taxi\textquotedblright{}.

\begin{cousinweb}
\textbf{\ar{تاكسي}} (\emph{tāksī}) is a modern loan from French \emph{taxi} (EtymDB); it has no Arabic root.
\end{cousinweb}

\begin{grammarlens}[title={What you've built}]
One more word for this chapter.
\end{grammarlens}

\subsection*{Your turn}
\begin{itemize}
\raggedright
  \item \emph{Say it:} \emph{tāksī}
  \item \emph{Say it:} \emph{tāksī}, once more
  \item \emph{Say it:} say \emph{tāksī}
  \item \emph{Recall:} say the Arabic for luggage, then the Arabic for a charger, then say \emph{tāksī} again
\end{itemize}

\begin{tcolorbox}[breakable,colback=teal!4,colframe=teal!35!black,title={Before you move on}]
What does \ar{تاكسي} mean? (\textquotedblleft{}A taxi\textquotedblright{}.) Say it once more.
\end{tcolorbox}
